\section{Progressive Reduced-Dimensional Demand Modelling}

The complete OD--time table is both computationally large and statistically
underdetermined by aggregate boarding and alighting counts.  The software
therefore resolves one composable model specification before JAX tracing.  A
disabled component creates no parameter block, design matrix, objective term,
or run-time branch.  This is a software design choice; the statistical
components themselves are established methods or explicit transit adaptations.

For origin-period group $g=(o,t)$ and cell $c=(o,d,t)$,
\[
 \log P_g=\log P_g^{(0)}+Z_g^{P}\theta_P,
 \qquad
 V_c=\log A_{dt}^{(0)}+Z_{dt}^{A}\theta_A-X_c\beta+I_c,
 \qquad
 q_c=P_g\frac{\exp V_c}{\sum_{k\in g}\exp V_k}.
\]
The fixed-routing mean remains $\boldsymbol\mu=B\boldsymbol q+\boldsymbol c$.
Period coefficients are the first relaxation because systematic temporal
residuals can otherwise be aliased into every spatial component.  Origin-group
effects modify productions, whereas destination-group effects modify relative
attractiveness; they are not interchangeable.

Automatic regions are deterministic nested partitions of normalized supply,
topology, and measurement-observability signatures.  They are model-resolution
regions, not socioeconomic zones.  The implementation adapts the automated
resolution and held-out-validation motivation of \cite{MenonEtAl2015} and the
homogeneity/contiguity hierarchy of \cite{YuanEtAl2019}; applying those ideas to
transit supply and observability is a methodological hypothesis requiring
empirical validation.

After margin heterogeneity has been assessed, a rank-$r$ interaction uses
\[
 I_{odt}=U_{G_O(o),:}V_{G_D(d),:}^{\mathsf T}.
\]
The columns of both factors are centered.  Ridge penalties stabilize their
scale, but rotational ambiguity means individual factor columns are not given
a behavioral interpretation.  Low rank follows the dimensional-reduction
insight used for time-varying transit OD inference by \cite{ChenChengSun2025};
the present JAX/MAP implementation is not a reproduction of that paper's MCMC
method.

Multiplicative gravity relationships are fitted in count scale rather than by
log-linear least squares, following the Poisson pseudo-likelihood argument of
\cite{SantosSilvaTenreyro2006}.  The negative-binomial parameterization has
mean $\mu$, dispersion $r$, and variance $\mu+\mu^2/r$
\cite{CameronTrivedi2013}.  These are the two count likelihood families
implemented by the current gravity objective.  Zero-inflated Poisson and
zero-inflated negative-binomial forms are useful future alternatives, but are
not available in the current gravity estimator and must not be presented as
production options.  In either case, a timetable structural zero remains
distinct from an observed zero: network-infeasible cells are removed before
estimation, whereas a zero-inflation component would describe retained
measurement records.

\begin{table}[htbp]
\centering\scriptsize
\setlength{\tabcolsep}{3pt}
\begin{tabular}{p{2cm}p{1.7cm}p{2.6cm}p{1.4cm}p{2.1cm}p{2.4cm}}
\hline
Component & Added parameters & Purpose & Cost & Risk & Diagnostic trigger\\
\hline
Period effects & $T-1$ per block & Temporal heterogeneity & Low & Reference/centering & Residuals by period\\
Origin groups & $G_O-1$ & Production heterogeneity & Low & Weakly observed origins & Origin-group residuals\\
Destination groups & $G_D-1$ & Attraction heterogeneity & Low & Attraction confounding & Destination residuals\\
Rank $r$ & $r(G_O+G_D)$ & Residual OD association & Moderate & Scale and rotation & OD pattern after margins\\
Negative binomial & 1 & Overdispersion & Low & Large-$r$ boundary & Variance beyond Poisson\\
ZIP/ZINB (future) & design columns & Separate zero process & -- & Not implemented in current gravity path & Excess zeros after NB\\
\hline
\end{tabular}
\caption{Optional components of the resolved demand family.}
\end{table}

\begin{table}[htbp]
\centering\small
\begin{tabular}{lll}
\hline
Stage & Mean structure & Observation family\\
\hline
M0 & Global production and impedance & Poisson\\
M1 & Add period coefficients & Poisson\\
M2 & M1 & Negative binomial\\
M3 & Add origin groups & Negative binomial\\
M4 & Add destination groups & Negative binomial\\
M5 & Add rank one & Negative binomial\\
M6 & Ranks 2, 4, then 8 & Negative binomial\\
Sensitivity & Smallest adequate mean & Future ZIP/ZINB only after implementation and validation\\
\hline
\end{tabular}
\caption{Recommended progressive ladder; advancement is diagnostic, not automatic.}
\end{table}

Warm starts copy common named blocks, expand period or hierarchy blocks under
their declared sum-to-zero convention, and initialize new interaction columns
with small deterministic values.  In-sample assignment residuals diagnose fit;
held-out measurements diagnose generalization and must never enter fitting.
Every block reports sensitivity, associated observations, gradient or
curvature evidence, penalty contribution, held-out improvement, and whether it
is data-supported, mixed, assumption-dominated, or unidentified.  No synthetic
``percentage from data'' is reported.  Baselines, grouping, gravity form,
regularization, low rank, structural-zero rules, fixed routing, zero inflation,
and identification constraints remain explicit assumptions.

The generic optimizer applies the same operational contract as the minimal
reduced-OD fitter.  Maximum likelihood and Gaussian-prior MAP are distinct
modes.  GravityModelSpecification transforms raw parameters to their physical
domain; the default SciPy gravity fit is unbounded in raw coordinates, while
the optional Biogeme pilot may receive explicit bounds.  An absolute time
budget requires an atomic checkpoint.  A stopped run records the latest
accepted iterate and resumes only when the complete model identity
(specification, layout, features, grouping, response operator, observations,
and software schema) matches.  Compilation time and warm value-and-gradient
time are reported separately from optimization time.
Each accepted iteration reports its duration and a bounded recent-duration
history (up to 32 iterations in the current gravity estimator), together with
the remaining iteration allowance, estimated remaining seconds, and expected
UTC completion time.  The estimate excludes JAX compilation and is therefore
meaningful only after measured iterations exist.  Progress callbacks are
optional at low level, while case-study drivers normally create a durable
progress sink; reporting adds no objective evaluations and is outside model
identity.  With a configured checkpoint, progress includes its path and age;
an interactive interrupt atomically saves the latest accepted iterate and
returns an explicit ``interrupted'' status suitable for later identity-
validated resume.

For migration of existing examples, a thin adapter exposes an established
fixed-routing gravity operator directly to the generic kernel.  If the legacy
OD rows require canonical reordering, the adapter permutes the demand vector on
device before the operator product; it never materializes a dense response
matrix.  Bounded end-to-end tests exercise this path on both synthetic public
networks and on the public Geneva GTFS snapshot.  On the latter (96 free OD
cells and 8,967 measurements, CPU, one optimization iteration), the generic M0
kernel compiled in 0.17 seconds and its warm value-and-gradient took
$1.1\times10^{-4}$ seconds.  These figures validate integration and kernel
cost, not statistical convergence; model comparison still requires the full
progressive fitting and held-out workflow described above.
